\section{Free external detection, quiet charges, and collinear tests}
\label{sec:physical-ward-free-classification}

This section proves three free-theory inputs used in the selected
finite-level Ward argument: a uniform finite external detector, the
classification of quiet vector charges, and reduction to finitely many
collinear tests.  Internal Ward-map erasure is treated separately in
Section~\ref{sec:free-internal-ward-erasure}.  The arguments here use
finite generation of a polynomial ideal, explicit Wick collision limits,
the strict all-grade free detector, and the signed symmetries of four coordinates.

\subsection{Polynomial grading and a finite external detector}

Write $X_{r,\alpha}$ for a polynomial coordinate of jet order
$r\in\mathbb N_0$; the finite color-slot indices $\alpha$ and the four
coordinate axes carry their fixed finite orderings.  Assign
weight $r+2$ to this coordinate.  Let $\mathcal R_j$ be the polynomial
subalgebra generated by coordinates of jet order at most $j$, and let
$\mathcal S_D$ be the gauge-invariant observables of weighted degree at
most $D$.  Denote by $\pi_k$ projection to ordinary polynomial length
$k$, counting each coordinate as one factor.  For $D,k\geq0$, define
$\mathcal H_{D,k}$ as the linear span of products
$v_1\cdots v_m$, where each $v_i\in\mathcal S_D$ is homogeneous of
ordinary length $k_i$ and $\sum_i k_i=k$; the empty product is allowed
when $k=0$.

An observable $p$ is \emph{null at capacity $D$} if
\begin{equation}
 \mathcal M(x;p,v_1,\ldots,v_m)=0
 \quad\text{for every }m,\ x\in\operatorname{Conf}_{m+1},
 \ v_i\in\mathcal S_D.
\label{eq:free-null-capacity}
\end{equation}
Here $\mathcal M$ is the free Euclidean moment and $p$ is the distinguished
first insertion.  This quantifies over arbitrary finite spectator words,
all separated positions, and all spectator observables in the fixed
capacity.

\begin{lemma}[Uniform generation at bounded jet order]
\label{lem:free-uniform-generation}
For each $j$ there is $D_j<\infty$ such that every gauge-invariant
$p\in\mathcal R_j$ homogeneous of ordinary length $k$ belongs to
$\mathcal H_{D_j,k}$, for every $k$.
\end{lemma}

\begin{proof}
There are finitely many coordinates of jet order at most $j$, so
$\mathcal R_j$ is a polynomial ring in finitely many variables and is
Noetherian.  Let $I$ be the ideal generated by all positive-length,
homogeneous, gauge-invariant elements of $\mathcal R_j$.  Choose finite
generators $g_1,\ldots,g_s$ of this ideal from that generating family,
and let $D_j$ bound their weighted degrees.  Each $g_i$ has positive
ordinary length $\ell_i$.

Induct on $k$.  Constants are in $\mathcal H_{D_j,0}$.  For positive
length, an invariant homogeneous $p$ is in $I$, so write it as a finite
sum $p=\sum_i r_i g_i$.  Let $\operatorname{Av}$ be compact-group
averaging.  Since $p$ and $g_i$ are invariant and averaging commutes with
length projection,
\[
 p=\sum_i \operatorname{Av}\bigl(\pi_{k-\ell_i}(r_i)\bigr)g_i,
\]
where terms with $\ell_i>k$ are zero.  Indeed, projecting $r_i g_i$ to
length $k$ gives exactly $\pi_{k-\ell_i}(r_i)g_i$, by additivity of
ordinary degree.  Each averaged first factor is invariant, homogeneous,
and of length $k-\ell_i<k$.  The induction hypothesis puts it in
$\mathcal H_{D_j,k-\ell_i}$, and multiplication by $g_i\in\mathcal S_{D_j}$
puts the term in $\mathcal H_{D_j,k}$.  This proves the claim with a
capacity independent of $k$.
\end{proof}

\begin{lemma}[Coalescence of bounded spectators]
\label{lem:free-coalescence}
If $p$ is null at capacity $D$ and $\pi_\ell p=0$ for every $\ell<k$,
then for any $q\in\mathcal H_{D,k}$ and any $x\ne y$,
\[
 \mathcal M((x,y);\pi_kp,q)=0.
\]
\end{lemma}

\begin{proof}
It is enough to take $q=v_1\cdots v_m$ with homogeneous factors of
lengths $k_i$ summing to $k$.  Choose distinct real numbers $c_i$ and
replace the spectator factors by insertions at $y+t c_i e_0$.  For all
sufficiently small positive $t$, these points are mutually distinct and
avoid $x$.  The null condition makes the resulting moment zero.  Expand
into Wick monomials.  By hypothesis, the distinguished components of
lengths below $k$ are zero and can be removed before taking any limit.  A
component of length above $k$ gives no Wick pairing at all: the distinguished
insertion has more than $k$ legs, while the entire spectator word has only
$k$ legs, and legs at the distinguished insertion cannot pair with one
another.  For a distinguished monomial of length exactly $k$, every
spectator leg must pair with a distinguished leg.  Hence there are no
spectator--spectator contractions, at any $t$.  The only kernels in these
surviving terms connect $x$ to a spectator point; they are smooth there
because $x\ne y$, and converge to their values at $y$ as $t\downarrow0$.
Their limit is precisely the two-insertion moment
$\mathcal M((x,y);\pi_kp,q)$.  Since the original moments vanish, this
limit is zero.  Extend by linearity to $\mathcal H_{D,k}$.
\end{proof}

\begin{lemma}[Shell reduction at a fixed jet cutoff]
\label{lem:free-shell-reduction}
Fix a jet cutoff $j$ and a gauge-invariant polynomial
$p\in\mathcal R_j$ null at capacity $D_j$.  Then $p=0$.
\end{lemma}

\begin{proof}
Induct on ordinary length $k$.  Suppose the components below $k$ vanish.
The coalescence lemma gives
\[
 \mathcal M((x,y);\pi_kp,q)=0
 \quad(q\in\mathcal H_{D_j,k},\ x\ne y).
\]
By Lemma~\ref{lem:free-uniform-generation}, the time-reflected conjugate
$\Theta\pi_kp$ belongs to $\mathcal H_{D_j,k}$: reflection preserves
the jet cutoff, the gauge-invariant sector, and ordinary length.  If
$\pi_kp\ne0$, the strict all-grade free detector
(Theorem~\ref{thm:ex-lattice-all-grade-free-detector}, with time
reflection and coefficient conjugation) at the separated pair
$(-r e_0,r e_0)$ gives
\[
 \mathcal M((-r e_0,r e_0);\pi_kp,\Theta\pi_kp)>0;
\]
reflection positivity alone would give only $\ge0$. This contradicts
the preceding vanishing.  Hence $\pi_kp=0$ for every $k$.
The polynomial is a finite sum of these components, so it is zero.
\end{proof}

To remove the fixed-jet restriction, put $j=\dim\mathfrak g$ and
$L=j+1$.  Suppose the gauge-invariant polynomial $p$ is supported in jet
orders at most $s+1$.  For each color $c$ choose a lower probe $q_c$ of
order $c$ and a sheet vector $w_c$ with $q_c(w_c)=1$.  These sheets are
nonzero: if the order-$c$ sheet were zero, every order-$c$ probe would be
zero, so the nonzero zero-order coordinate cell (which takes value one on
its matching basis two-form) would pair to zero with all such probes at
every nonzero displacement, contrary to the all-order separation argument
below.  The replacement equations have polynomial
matrix entries
\[
 M_{ci}=m_{ci},
\]
where $m_{ci}$ is the coordinate polynomial obtained by evaluating the
order-$c$ probe $q_c$ on the field coordinate of color $i$ at jet order
$c$; $i$ ranges over the color slots of the order-$(s+1)$ derivative.  Its
determinant is not the zero polynomial.  To see this, evaluate the Taylor
tower by putting $w_c$ in the matching slot of color $c$ and zero in every
other color slot.  The matrix becomes the identity.  This specialization
proves that its determinant polynomial is nonzero; it does not claim that
the matrix is invertible at every field value.

Fix any upper probe of order $L$ and any nonzero displacement.  Let
$R_i$ be the finite linear combination of order-$(s+1)$ partial
derivatives whose coefficients are the free two-point pairings of that
upper probe with the color-$i$ jet coordinates.  The circulation
equations for the chosen $q_c$ state that $M R=0$ as a polynomial
identity.  Multiplying by the adjugate gives the determinant times each
component equal to zero.  The polynomial coordinate ring is an integral
domain, so each component $R_i$ vanishes.  This holds for every upper
probe and every nonzero displacement.  We use the following separation
fact at all jet orders: if a jet functional of order $h$ pairs to zero at
every nonzero displacement with every probe of order $\ell$, it is zero.
Induct on $\ell$.  At $\ell=0$, its color-resolved pulse vanishes off the
insertion by the free-kernel formula.  Testing the single-charge radial
flux identity~\eqref{eq:maxwell-radial-single-charge} against arbitrary
homogeneous Maxwell waves of order $h$ gives $-(h+2)$ times its value on
each order-$h$ jet, hence all coefficients vanish.  For $\ell>0$, fix a
probe $q$ of order $\ell-1$ and let $K(x)$ be its pairing with the order-$h$
functional.  Every directional derivative of $K$ is a pairing against a
raised probe of order $\ell$, so it vanishes by hypothesis.  The kernel's
homogeneity gives the Euler identity
$x\cdot\nabla K(x)=-(4+h+\ell-1)K(x)$.  Its coefficient is nonzero,
so $K$ vanishes at every nonzero $x$; the induction hypothesis now applies.
Therefore the color-resolved two-point functionals,
indexed by every nonzero displacement and every upper probe of the chosen
order, span the full dual of the order-$(s+1)$ jet space.  Applying this
spanning statement to
each color slot shows that every ordinary partial derivative in an
order-$(s+1)$ variable vanishes.  In characteristic zero, $p$ is
independent of those variables and lies in the lower jet shelf.

Suppose the gauge-invariant polynomial $p$ is null at capacity $D_*$ and
lies in weighted shell
$w+1$.  Whenever the highest jet order satisfies $s+1>j$, test the circulation
with each lower probe of order $\ell<j$; then $\ell<s+1$, as required
for weighted-degree lowering.  The upper probe has order $L=j+1$, so the
separate collision condition $\ell<L$ also holds.  The mixed-collision nullity theorem
(Section~\ref{sec:free-collision-pole-separation},
Theorem~\ref{thm:free-circulation-null}) now applies: the current
polynomial has jet cutoff $s+1$ and is a ghost at capacity $D_*$; the
circulation uses upper probe order $k=L$, lower order $\ell$, nonzero
displacement, and satisfies $\ell<L$ and $L+\ell+4\le D_*$.  It
therefore remains a ghost at capacity $D_*$.  The admissibility bound
$L+\ell+4\leq D_*$, with $L=j+1$ and $\ell<j$, is exactly what allows
this upper and lower probe in the extra test.  Each term replaces a coordinate of weight $s+3$ by
one of weight $\ell+2$, so the circulation lies in weighted shell $w$.
Strong induction on $w$ therefore makes it zero.  The top-shelf
elimination then applies.

Iterating lowers the jet shelf to the fixed order $j$.  If $r_j$ is the
host bound from finite generation, take
\[
 D_* = \max\{r_j,\,4+j+(j+1)\}.
\]
The second term dominates $(j+1)+\ell+4$ for all lower probe orders
$\ell<j$.  The fixed-shelf argument above then annihilates the polynomial.  This proves the following uniform result.

\begin{theorem}[Uniform external detection]
\label{thm:free-uniform-detection}
For a compact simple gauge group there is a finite $D_*$ such that every
gauge-invariant observable in every finite shell which is null at capacity
$D_*$ is zero.
\end{theorem}

\subsection{Quiet charges and collinear tests}

For a four-component vector of linear operators $A=(A_a)_{a=1}^4$ on the
finite-level field space, say that $A$ is quiet at capacity $D$ if, for
each component $a$ and every separated word with labels in $\mathcal S_D$,
the sum of moments obtained by applying $A_a$ to the label in each slot
and summing over the slots vanishes.  Write
$\partial_a$ for the canonical translation derivative, raising jet order
by one and extended by the Leibniz rule.

\begin{theorem}[Classification of quiet charges]
\label{thm:free-quiet-classification}
Let $\mathcal P_k$ be the finite field space at level $k$, and let
$A=(A_a)_a$ be a four-component charge in the specified vector-spin
subspace.  Thus $A_a(\mathcal P_k)\subseteq\mathcal P_{k+1}$ for every
$k\leq d$, the maps have the signed-coordinate covariance and reality
conditions of the vector representation, and they annihilate the unit.  If $d\geq\max\{4,D_*\}$ and the charge is quiet at
capacity $d$, then there is one real scalar $c$ with
$A_a=c\partial_a$ for all four axes.
\end{theorem}

\begin{proof}
Let $S_0$ and $S_1$ be the real sheet spaces of zero-jet and one-jet
curvature coordinates, and put $V_i=S_i^*$.  If $p$ is a linear functional
on a sheet, write $\phi_{r,c}(p)$ for its color-$c$ field coordinate and
let $\mathsf Q(p,q)$ denote the color-contracted quadratic
$\sum_c\phi_{r,c}(p)\phi_{s,c}(q)$, with the jet orders indicated by the
spaces of $p,q$.  The degree-four shell is the direct sum of the unit and
the span of
$\mathsf Q(p,q)$ for $p,q\in V_0$.  The degree-five shell is the sum
of that degree-four shell and the mixed products $\mathsf Q(p,q_1)$ with
$p\in V_0$, $q_1\in V_1$.  The charge kills the unit.  By the vector
filtration and opposite-frame parity, its response to a degree-four
quadratic stress label lies in the mixed degree-five subspace.  It is real
because the charge preserves the real structure.  Each mixed product is
real; taking the real part of a complex linear combination shows that a
real element of their complex span lies in their real span.  The mixed products therefore represent that response by a real linear
panel $D_{a,p}:S_0\to V_1$: in a basis $(e_\alpha)$ of $S_0$ with dual
basis $(e^\alpha)$, its mixed polynomial is
$\sum_\alpha\mathsf Q(e^\alpha,D_{a,p}(e_\alpha))$.  This expression
is basis-independent because it is the contraction of a linear map with
the identity tensor.  The representation is unique: evaluate the
polynomial on a Taylor tower with independently prescribed order-zero
and order-one sheets in a fixed common color slot and all other slots
zero.  The value is exactly $D_{a,p}(X)(Y)$, so vanishing of the
polynomial forces every panel coefficient to vanish.  This constructs the
finite-level panel without assuming any formula for the charge on higher
shells.

Write $\mathcal M_3(x,y,z;A,B,C)$ for the free moment with labels
$A,B,C$ inserted at $x,y,z$.  The free six-leg Wick computation used here
is the following identity:
\begin{equation}
\begin{split}
&\mathcal M_3(x,y,z;\mathsf Q(a,v),\mathsf Q(p,p),\mathsf Q(q,q))\\
&\quad=4\dim(\mathfrak g)\bigl(F(a,p;x-y)F(v,q;x-z)
 +F(a,q;x-z)F(v,p;x-y)\bigr)F(p,q;y-z).
\end{split}
\label{eq:free-three-quadratic-wick}
\end{equation}
It follows by pairing the two legs at the first insertion with the two
legs at the other insertions in the two possible assignments; the last
pair joins the remaining sites.  The color deltas leave one free color
sum, and the two choices of the leg within each quadratic give the factor
$4$.  For a general panel, expand in a basis of $S_0$ and use
multilinearity.  Choose the third quadratic with label $w$ so that
$F(q,w;-z)=q(X)$ for every $q\in V_0$.  To see this explicitly, write
$R_z$ for the invertible exterior-square action of the House reflection
associated with $z$ (so $R_{-z}=R_z$), and put
$\kappa_z=4\nu\lVert z\rVert^{-4}>0$, where $\nu=(4\pi^2)^{-1}$ is the positive
free-kernel normalization.  The zero-jet covariance is
\[
 F(q,\operatorname{cell}_{ij};z)=\kappa_z\,q(R_z\operatorname{blade}_{ij}).
\]
For given $X$, set $W=R_z^{-1}X$ and define
\[
 w=\kappa_z^{-1}\sum_{i,j}\frac12
   \operatorname{cell}_{ij}(W)\,\operatorname{cell}_{ij}.
\]
The coordinate resolution
$W=\sum_{i,j}\frac12\operatorname{cell}_{ij}(W)\operatorname{blade}_{ij}$
then gives $F(q,w;z)=q(X)$ for every $q$.  Thus this third quadratic
exists for each $X\in S_0$ and nonzero $z$.

Fix a vector component $a$ and abbreviate $D_p=D_{a,p}$.  Insert the
three quadratic probes at $tu,0,z$, multiply their moment by
$t^5$, and let $t\downarrow0$.  The identity above gives the first
near-collision limit.  Exchanging the first two insertion points gives a
second near-collision limit with $p,q$ interchanged and displacement
$-u$; the one-jet/zero-jet kernel is odd, so this term has the opposite
sign.  The third placement has its first insertion at the fixed point $z$
and its scaled limit is zero.  Every separated three-point Ward row
vanishes, so the sum of these limits is zero.  Cancelling the common
nonzero factor $4\dim(\mathfrak g)$ gives, for every $X\in S_0$ and every
nonzero separation $u$,
\[
 F(D_p(X),q;u)\,q(X)=F(D_q(X),p;u)\,p(X),
\]
where $F(v,q;u)$ is the free pairing for $v\in V_1$, $q\in V_0$ at
separation $u\ne0$.  If $p(X)=0$, the identity gives
$F(D_{a,p}(X),q;u)=0$ for every $q$ with $q(X)\ne0$ and every nonzero
$u$.  For $X\ne0$, these $q$ form the complement of a hyperplane and
span $V_0$, so linearity extends the equality to every upper probe.  For
$X=0$ the panel already vanishes by linearity.  We justify the separation
used here and below.  If $v\in V_1$ pairs to zero with every zero-jet
probe at every nonzero displacement, the free-kernel formula shows that
the color-resolved pulse of $v$ vanishes componentwise away from its
marked point.  Test the single-charge radial-flux
identity~\eqref{eq:maxwell-radial-single-charge} against an arbitrary
order-one sheet $A\in S_1$, represented by its homogeneous Maxwell wave,
using a compact cutoff equal to one near the marked point.
The flux vanishes because the pulse vanishes on
the support of the cutoff gradient.  Since $v$ has jet order one, the
radial identity evaluates the same flux as $-3$ times the pairing of $v$
with $A$ at the marked point.  Since $A$ is arbitrary in $S_1$, this gives
$v=0$.  Thus the
free pairings over all nonzero displacements have zero annihilator and span
$V_1^*$.  In particular $D_{a,p}(X)=0$ on $\ker p$.  For $p\ne0$, choose
$X_0$ with $p(X_0)=1$ and define $f_a(p)=D_{a,p}(X_0)$.  Linearity of the
panel then gives $D_{a,p}(X)=p(X)f_a(p)$ for every $X$.  Set $f_a(0)=0$.
Thus $D_{a,p}=p\otimes f_a(p)$ and
$A_a(\mathsf Q(p,p))=\mathsf Q(p,f_a(p))$.  The Taylor-tower
extraction above shows this panel representation is unique.  For nonzero
$p,q$, choose $X$ with $p(X)q(X)\ne0$ and cancel these scalar factors in
the triangular identity; when either probe is zero, the resulting
identity is immediate.  The coefficient
identity becomes the reciprocity relation
\[
 F(f_a(p),q;z)=F(f_a(q),p;z),\qquad z\ne0.
\]
For fixed $a$, reciprocity and bilinearity of $F$ imply additivity:
for every $r,z$,
\[
 F(f_a(p+q)-f_a(p)-f_a(q),r;z)=0.
\]
The same computation with $p$ replaced by $tp$ gives real homogeneity.
The free two-point functionals over all nonzero displacements separate
$V_1$, so $f_a$ is real-linear.
Signed-coordinate covariance gives
\[
 \rho_s f_a(p)=\varepsilon(s,a)f_{s a}(\rho_s p),
\]
with $\rho_s$ the signed action and $\varepsilon(s,a)$ its coordinate sign.

Here is the finite coefficient determination.  Let $C_{ij}$ be the six
coordinate two-forms in $V_0$ and let $D_{a;ij}$ be their one-jet curls in
$V_1$.  Signed-reflection parity allows only
$D_{0;01},D_{2;12},D_{3;13}$ in $f_0(C_{01})$.  Equivariance under the
swap of axes $2,3$ makes this vector invariant under that swap, so it lies
in the span of $D_{0;01}$ and $D_{2;12}+D_{3;13}$.  The free-shift
identity (the divergence identity)
\[
 D_{2;12}+D_{3;13}=D_{0;01}
\]
therefore gives $f_0(C_{01})=A D_{0;01}$.  For $C_{12}$,
signed-reflection parity first allows exactly
\[
 D_{0;12},\quad D_{1;02},\quad D_{2;01}.
\]  The negative eigenspace of the swap of
axes $1,2$ then makes the last two occur as their difference.  The Bianchi identity
\[
 D_{1;02}-D_{2;01}=D_{0;12}
\]
then gives
$f_0(C_{12})=B D_{0;12}$.  Evaluate the reciprocity relation at the
nonzero cross-seed $z=e_0+e_2$.  By the explicit cross-seed evaluation
$F(D_{0;01},C_{12};e_0+e_2)=-2\,\nu$, with
$\nu>0$, the common factor is nonzero and $A=B$.  Signed permutations fixing axis $0$ then carry these coefficient
identities to the remaining coordinate two-forms.  Since the $C_{ij}$ form
a basis, this proves $f_0=A\nabla_0$, where $\nabla_0$ is the
zero-to-one-jet derivative insertion; covariance gives
$f_a=A\nabla_a$ for every axis.  The Leibniz rule gives
\[
 \partial_a\mathsf Q(p,p)=2\mathsf Q(p,\nabla_a p),
\]
so the original charge is $c\partial_a$ on the whole degree-four shell,
with $c=A/2$.  Its unit rule supplies the constant component.  Subtract
$c\partial_a$ and call the remainder $B_a$.  This remainder is quiet and
vanishes on the unit and on every quadratic
stress insertion.  Let $\Gamma$ be the grading derivation that assigns
weight $r+2$ to a coordinate of jet order $r$.  The localized flux identity
gives for each $p$ in the finite shell that
\[
 \Gamma(B_a p)-B_a(\Gamma p)
\]
has zero Ward rows at the same external capacity.  The required dilation Ward identity is the radial flux formula of
Section~\ref{sec:free-maxwell-radial-flux}, equation~\eqref{eq:maxwell-radial-polynomial-ward}.  For a smooth cutoff
$\chi$ equal to one near the marked point $x$, the centered dilation
field is $v(y)=y-x$.  On a homogeneous Maxwell jet of order $r$, its
infinitesimal action is $(y-x)\cdot\nabla+2$; Euler's identity gives
$(r+2)$ times that jet at $x$.  The radial flux formula identifies the
cutoff stress integral with the negative of this insertion.  Expanding a
polynomial into its finite Wick monomials and summing these single-jet
identities gives exactly the weighted Euler operator $\Gamma$ on the
marked polynomial.  Choose $\chi$ equal to one at the marked insertion
and zero at the other insertions.  The ordinary Ward terms vanish because
$B_a$ is quiet; the boundary term vanishes because $B_a$ annihilates the
quadratic stress labels.  Therefore the commutator has zero rows against
all separated test words at capacity $d$.  Since $d\geq D_*$, it is null
at capacity $D_*$.  Theorem~\ref{thm:free-uniform-detection}
therefore gives $\Gamma B_a=B_a\Gamma$ on the finite shell.

On a monomial $X^\nu$, $\Gamma$ has eigenvalue
$m=\sum_{r,\alpha}(r+2)\nu_{r,\alpha}$.  Opposite-coordinate reflection
has eigenvalue $(-1)^{\sum r\nu_{r,\alpha}}$.  These signs agree because
the exponents differ by twice the ordinary length.  Spin covariance makes
$B_a$ anticommute with this reflection.  Since commutation with $\Gamma$
preserves each weight eigenspace and reflection is scalar $(-1)^m$ on that
eigenspace, anticommutation forces $B_a$ to vanish there.  Every polynomial
is a finite sum of weight-homogeneous components, so $B_a=0$.  The initial
coefficient comparison gave one common scalar across the four axes; this
proves the classification.
\end{proof}

\begin{theorem}[Finite collinear scene family]
\label{thm:free-collinear-scenes}
\begin{enumerate}
\item At every level $d$ there is a finite family of collinear scene
  tests whose simultaneous kernel on the space of charges is exactly the
  quiet-charge subspace.
\item If $d\geq\max\{4,D_*\}$, the simultaneous kernel of such a family on
  the vector-spin subspace of Theorem~\ref{thm:free-quiet-classification}
  is exactly $\mathbb R\partial$.
\end{enumerate}
\end{theorem}

\begin{proof}
Let $V$ be the finite-dimensional real space of charges at the chosen
level.  For a charge component $A_a$, a scene test is the complex-valued
real-linear functional obtained by evaluating an arbitrary finite
configuration derivative of its full Ward sum at a strict collinear
configuration, using a line axis different from $a$.  The intersection of the kernels of all such tests is the
quiet subspace.  One direction is immediate.  Conversely, fix a separated
configuration $x$ and a charge component $a$.  Choose a line axis $b\ne a$.
Order the indices so that all strict
orders of the projections $x_i\cdot e_b$ are preserved, and choose an
arbitrary order within each tie.  For sufficiently small $\delta>0$, the
perturbed projections
$q_i=x_i\cdot e_b+\delta\,\mathrm{rank}(i)$ are pairwise distinct and
preserve every previous strict inequality.  Put $y_i=q_i e_b$.  For each
pair, its projected difference along the segment from $y$ to $x$ keeps the
same strict sign until the endpoint; at the endpoint the original
configuration is separated.  Thus the segment stays in configuration
space.
The moment is analytic on
the separated configuration space.  If all derivatives at the collinear
endpoint vanish, its Taylor series vanishes near that endpoint, and the
one-variable identity theorem propagates vanishing along the segment.
Thus all separated Ward rows vanish, which is quietness.

An arbitrary intersection of kernels on a finite-dimensional space is
already the intersection of a finite subfamily: add a functional whenever
the current intersection is too large.  Its dimension strictly drops,
so after at most $\dim V$ additions the process stops.  The selected
finite scene family has the required kernel, proving~(1); no level
threshold is used.

For~(2), intersect with the vector-spin subspace. By~(1) the kernel there
consists of the quiet charges in that subspace, and for
$d\geq\max\{4,D_*\}$ Theorem~\ref{thm:free-quiet-classification} shows
that each of them is $c\partial$ for a real $c$. Conversely $\partial$
lies in the vector-spin subspace (it raises the level by one, has the
signed-coordinate covariance and reality of a vector, and annihilates
the unit) and is quiet: applying $\partial_a$ in each slot and summing
gives the derivative of the free moment along a common translation of
all insertion points, which vanishes by translation invariance. Hence the
kernel on the vector-spin subspace is $\mathbb R\partial$.
\end{proof}

The finite-scene theorem is the free-theory interface used by the selected finite-level Ward package; its formal source map and remaining hypothesis checks are in the adjacent evidence record.
